
\PassOptionsToPackage{full}{textcomp}
\documentclass[]{tufte-handout}
%\usepackage{fontspec}

% load babel package and options here
%\usepackage[p,osf]{ETbb} % osf in text, tabular lining figures in math
\usepackage{ETbb} % osf in text, tabular lining figures in math
\usepackage[scaled=.95,type1]{cabin} % sans serif in style of Gill Sans
\usepackage[varqu,varl]{zi4}% inconsolata typewriter
\usepackage[T1]{fontenc} % LY1 also works
\usepackage[libertine,vvarbb]{newtxmath}
%\usepackage[cal=boondoxo,bb=boondox,frak=boondox]{mathalfa}

%\geometry{showframe} % display margins for debugging page layout

\usepackage{graphicx} % allow embedded images
%  \setkeys{Gin}{width=\linewidth,totalheight=\textheight,keepaspectratio}
  \graphicspath{{graphics/}} % set of paths to search for images
\usepackage{amsmath}  % extended mathematics
\usepackage{booktabs} % book-quality tables
%\usepackage{units}    % non-stacked fractions and better unit spacing
\usepackage{multicol} % multiple column layout facilities
\usepackage{lipsum}   % filler text
\usepackage{fancyvrb} % extended verbatim environments
  \fvset{fontsize=\normalsize}% default font size for fancy-verbatim environments
\usepackage{gensymb} % provides symbols like \degree
\usepackage{ragged2e} % enables hyphenation in ragged-right justification
\usepackage[normalize-symbols]{textalpha} %enables \textalpha for alpha symbol etc.

\usepackage{hyperref} % enables styling of href and url
\hypersetup{
    pdftitle={Mechanisms},
    pdfauthor={Barry Linkletter},
    colorlinks=true,
    linkcolor=blue,
    filecolor=magenta,      
    urlcolor=blue,
    pdfborder={0 0 0},
    frenchlinks=false,
    pdfpagemode=FullScreen,
    }
    \urlstyle{same}
\makeatletter
% Inspired by http://anti.teamidiot.de/nei/2009/09/latex_url_slash_spacingkerning/
% but slightly less kern and shorter underscore
\let\UrlSpecialsOld\UrlSpecials
\def\UrlSpecials{\UrlSpecialsOld\do\/{\Url@slash}\do\_{\Url@underscore}}%
\def\Url@slash{\@ifnextchar/{\kern-.03em\mathchar47\kern-.15em}%
    {\kern-.0em\mathchar47\kern-.08em\penalty\UrlBigBreakPenalty}}
\def\Url@underscore{\nfss@text{\leavevmode \kern.06em\vbox{\hrule\@width.3em}}}
\makeatother

\usepackage{enumitem} % allows resuming enumerate lists.
\usepackage{mathtools}
\usepackage{mhchem}

\usepackage{siunitx} % provides "S" column class for aligning decimals.  
         \sisetup{uncertainty-mode = separate}

\usepackage{nicefrac}

\usepackage[nospace]{varioref}
    \renewcommand\reftextfaceafter{on the following page}
    \renewcommand\reftextafter {on the next page}
    \renewcommand\reftextfacebefore{on the previous page}
    \renewcommand\reftextbefore {on the previous page}

\newcommand{\tss}[1]{\textsuperscript{#1}}

\usepackage{babel}
\usepackage{float}
\usepackage{stackrel}

\usepackage[shortconst]{physconst}

\usepackage[normalem]{ulem}  % provides strikethrough \sout{}

\usepackage{newfloat}
\DeclareFloatingEnvironment[
  fileext = los ,
  listname = {List of Schemes} ,
  name = Scheme
]{scheme}

\usepackage{newfloat}
\DeclareFloatingEnvironment[
  fileext = loe ,
  listname = {List of Eqs} ,
  name = Equation
]{equate}

\usepackage{listings}

%%%%%%%%%%%%%%%%%%%%%%%%%%%%%%%%%%%%%%%%%%%%%%%%%%%%%%%%%%%   


\title{Class Meeting \#15: Pre-Meeting Activity}
\author[Barry Linkletter]{}
\date{} % without \date command, current date is supplied
%%%%%%%%%%%%%%%%%%%%%%%%%%%%%%%%%%%%%%%%%
% From Williams Ch. 2 problem 2
%%%%%%%%%%%%%%%%%%%%%%%%%%%%%%%%%%%%%%%%%   
   


\begin{document}
\justifying


\maketitle% this prints the handout title, author, and date

\begin{abstract}

%\marginnote[-20mm]{\underline{\allcaps{your name here}}}
\noindent
Why don't we see \textit{ortho} subtituents in Hammett's substituent parameters? Let's graph some data and consider the results. This is an assigned problem to be completed before meeting~\#15. Bring a printed copy of your plot. 
\marginnote[-10mm]{This document was produced using the \LaTeX\ typesetting language with the Tufte-handout document class. Chemical diagrams were prepared in \emph{ChemDoodle}.}
\end{abstract}
%\printclassoptions


\section{Including \textit{ortho} Substituents}\label{sec:plan}


It is easy to measure the electronic effect of \textit{ortho} substituents on the acid equilibrium of benzoic acids and from that data calculate a set of $\sigma$ values for these substituents. If you examine the Williams data set of substituent parameters\sidenote
[][-15mm]{See the \texttt{LFER-Williams.csv} file available in the data repository for the course at \url{https://raw.githubusercontent.com/blinkletter/LFER-QSAR/refs/heads/main/data/LFER_Williams.csv}.\vspace{2mm}\label{ref:data}}
you will see that there are $\sigma$ values for a collection of \textit{ortho} substituents present.

In table~\vref{tab:rates2}, the rates for base hydrolysis of various ethyl benz\-oates\sidenote[][-0mm]{"Kinetic Studies of the Hydrolysis of Esters. I. Alkaline Hydrolysis of Esters of Substituted Benzoic Acids." E. Tommila, \textit{Ann. Acad. Sci. Fennicae}, \textbf{1941}, \text{Ser. A 57}, 3.\\ \vspace{2mm}\label{ref:tommila}}\tss{,}\sidenote[][40mm]{I could not find the Tommila paper, but our librarian did. I then obtained a copy through interlibrary loan.  \\ \vspace{2mm} This data is available in other collected works and reviews throughout the literature. For example, you can find this data in ``Free Energy Relationships in Organic and Bio-organic Chemistry'' by Andrew Williams, \textit{The Royal Society of Chemistry, Cambridge}, \textbf{2003}, pp 46. \url{https://doi.org/10.1039/9781847550927}. We have \href{https://ebookcentral.proquest.com/lib/upei/detail.action?docID=1185543}{access to the E-textbook version} via the UPEI library. \\ \vspace{2mm}

Williams reported the values calculated from Arrhenius plots of the kinetics and not the experimental value at \qty{25}{\degreeCelsius}. This was a way to include the data at other temperatures in the final reported value. Using the raw data from the paper, we could construct an Eyring plot for each result in table~\ref{tab:rates2} and then calculate our own values. We would also be able to calculate the standard error for the results. Would we observe an isokinetic temperature? Would the errors be significant? This could be the start of an exploration\ldots \label{ref:williams} \\ \vspace{2mm}

When you're serious, always seek out the  source. If you want to read the original paper, just ask me for a copy.} are presented.  Construct a Hammett plot and comment on the quality of the plot when the \textit{ortho} substituents are included and when they are not. Bring a copy of your plot to the class meeting for show-and-tell.

%The above publication is unavailable. the data in table~1 is actually from ``Free Energy Relationships in Organic and Bio-organic Chemistry'' by Andrew Williams, \textit{The Royal Society of Chemistry, Cambridge}, \textbf{2003}, pp 46. \url{https://doi.org/10.1039/9781847550927}. We have \href{https://ebookcentral.proquest.com/lib/upei/detail.action?docID=1185543}{access to the E-textbook version} via the UPEI library.\\\vspace{2mm}\noindent \textsc{Update}:

%I can now confirm that the data reported in Williams is the same. Williams reported the values calculated from Arrhenius plots of the kinetics and not the experimental value at \qty{25}{\degreeCelsius}. This was a way to include the data at other temperatures in the final reported value. Using the raw data from the paper, we could construct an Eyring plot for each result in table~\ref{tab:rates2} and then calculate our own values. We would also be able to calculate the standard error for the results. Would we observe an isokinetic temperature? Would the errors be significant? Someday we should do this. \label{ref:williams}} Construct a Hammett plot and comment on the quality of the plot when the \textit{ortho} substituents are included and when they are not. Bring a copy of your plot to the class meeting for show-and-tell.

%\begin{margintable}[-0mm]
\begin{table}[h!]
    \small
    \centering
    \fontfamily{ppl}\selectfont
    \caption[]{Rates of base hydrolysis for various ethyl benzoates at  \qty{25}{\degreeCelsius}.\tss{\ref{ref:tommila}}\tss{,}\tss{\ref{ref:williams}}\\ \vspace{2mm} 
    \noindent  You can download this data set from \url{https://raw.githubusercontent.com/blinkletter/4410PythonNotebooks/refs/heads/main/Class_17/data/4410_15-PreClassActivityData.csv}.
    }
    \vspace{2mm}
    \begin{tabular}{lSS}
      %  \toprule
{Substituent}      &  {$\sigma$}  &      {$k_{OH}$ /\num{e-3} \unit{M^{-1}s^{-1}}} \\
        \midrule
\ce{m-OCH3}       &           0.12 &      3.92         \\ 
\ce{p-N(CH3)2}    &          -0.83 &      0.0634       \\ 
\ce{p-Cl}         &           0.23 &     11.7          \\ 
\ce{m-Cl}         &           0.37 &     18.2          \\ 
\ce{p-Br}         &           0.23 &     13.9          \\ 
\ce{m-Br}         &           0.39 &     17.9          \\ 
\ce{p-I}          &           0.18 &     12.2          \\ 
\ce{m-I}          &           0.35 &     15.0          \\ 
\ce{p-F}          &           0.06 &      5.86         \\ 
\ce{p-CN}         &           0.66 &    157            \\ 
\ce{m-CN}         &           0.56 &    103            \\ 
\ce{m-NO2}        &           0.72 &    137            \\ 
\ce{p-NO2}        &           0.78 &    246            \\ 
\ce{p-NH2}        &          -0.66 &      0.0864       \\ 
\ce{m-NH2}        &          -0.16 &      1.66         \\ 
\ce{H}            &           0    &      2.89         \\ 
\ce{m-CH3}        &          -0.07 &      1.69         \\ 
\ce{p-CH3}        &          -0.17 &      1.14         \\ 
\ce{o-CN}         &           1.06 &    122            \\ 
\ce{o-OEt}        &          -0.01 &      1.16         \\ 
\ce{o-CH3}        &           0.29 &      0.338        \\ 
\ce{o-Cl}         &           1.26 &      4.40         \\ 
\ce{o-NO2}        &           2.03 &     16.9          \\ 
\ce{o-Br}         &           1.35 &      3.00         \\ 
\ce{o-I}          &           1.34 &      1.63         \\ 
    \end{tabular}
    \label{tab:rates2}
\end{table}
%\end{margintable}


\nobibliography{}

\end{document}